\scanpage{171}
\mspage{(20)}
SYSTÈMES de RACINES\hfill Tribu n$^{\handwritten{o}}$ \unsure{\handwritten{59}|59}

de signe $+1$.

1) On pose $E=\unsure{\mathbb{I}|\Gamma|\mathbb{P}|L|\Pi}\handwritten{/}\unsure{\handwritten{3}|3}P$. C'est un $\uuline{F}_{\handwritten{3}}$-espace vectoriel de dimension \unsure{5|\handwritten{5}}. Montrer que le produit $(x,y)$ définit sur $E$ une forme bil. sym. non dégénérée. Montrer que deux éléments distincts de $R$ ont des images distinctes dans $E$.

2) On note $\operatorname{O}(E)$ le groupe orthogonal de $E$, muni de la forme bil. précédente. On a $\operatorname{O}(E)=\handwritten{\{\pm}1\handwritten{\}}\times\operatorname{SO}(E)$. La norme spinorielle définit un homomorphisme surjectif $\operatorname{SO}(E)\handwritten{\longrightarrow\{\pm}1\handwritten{\}}$ dont le noyau est noté $\operatorname{SO}^{\handwritten{+}}(E)$. Le groupe $\operatorname{SO}^{\handwritten{+}}(E)$ est simple, (d'ordre \unsure{25 920|\handwritten{25 920}}) ; le quotient $\operatorname{O}(E)/\operatorname{SO}^{\handwritten{+}}(E)$ est de type $(2,2)$. En déduire que $\operatorname{O}(E)$ contient un sous-groupe $\handwritten{\Omega}(E)$ d'indice 2 distinct de $\operatorname{SO}(E)$ et ne contenant pas $-1$ ; montrer que ce sous-groupe est unique.

3) Tout élément de $\operatorname{Aut}(R)$ définit par passage au quotient un élément de $\operatorname{O}(E)$. Montrer, (en comparant les ordres), que l'on obtient ainsi un isomorphisme de $\operatorname{Aut}(R)$ sur $\operatorname{O}(E)$ ; l'image de $W(R)$ par cet isomorphisme est $\handwritten{\Omega}(E)$, celle de $W^{\handwritten{+}}(R)$ est $\operatorname{SO}^{\handwritten{+}}(E)$.

(Ainsi, $W(R)$ est extension d'un groupe à deux éléments \unsure{oar|par} le groupe simple d'ordre 25 920).

F) Groupe de Weyl $W$ de $E_{\unsure{\handwritten{7}|7}}$.

Soient $Q$ le réseau radiciel, $P$ celui des poids. On a $2P\handwritten{\subset}Q$. Posons $E=Q\handwritten{/}2P$. La forme bilinéaire $(x\ y)$ définit, par passage au quotient, une forme bilinéaire sur le $F_{\unsure{\handwritten{2}|2}}$-espace vectoriel $E$, (noter que $(Q\mid 2P)\handwritten{\subset}2Z$). Montrer que cette forme est non dégénérée et alternée, (noter que $\dim(E)=\unsure{\handwritten{6}|6}$). En déduire un homomorphisme $W\handwritten{\longrightarrow}\operatorname{Sp}(E)=\operatorname{Sp}(3,F_{\handwritten{2}})$.

Montrer que la suite
\[
\text{(s)}\qquad 1\handwritten{\longrightarrow}(1,-1)\handwritten{\longrightarrow}W\handwritten{\longrightarrow}\operatorname{Sp}(3,F_{\handwritten{2}})\handwritten{\longrightarrow}1
\]
est exacte, (on peut utiliser le fait que $\operatorname{Sp}(3,F_{\handwritten{2}})$ est d'ordre $2^{\unsure{\handwritten{9}|9}}.3^{\unsure{\handwritten{4}|4}}.5.7$ ; cf. Alg. IX).

Montrer que le sous-groupe $W^{\handwritten{+}}$ formé des éléments de déterminant 1 de $W$ se projette isomorphiquement sur $\operatorname{Sp}(3,F_{\handwritten{2}})$, (noter que ce sous-groupe ne contient pas $-1$).

Donner une seconde démonstration de (s) en utilisant la structure de $W(E_{\unsure{\handwritten{8}|8}})$ et l'isomorphisme $\operatorname{Sp}(3,F_{\handwritten{2}})\handwritten{\longrightarrow}\uline{\operatorname{O}}(7,F_{\handwritten{2}})$.
% archive sheet number: 170
